\documentclass[12pt]{article}
\usepackage{amsmath, amssymb}
\usepackage[margin=1in]{geometry}
\setlength{\parskip}{6pt}
\title{QUBO Formulation: Remote-Sensing Image Segmentation}
\date{}
\begin{document}
\maketitle

\subsection*{Remote-Sensing Image Segmentation}

\paragraph{Problem Statement.}
Given a remote-sensing image represented as a grid graph $G = (V, E)$, where $V$ denotes the set of all pixels and $E$ denotes the set of adjacency relationships between neighboring pixels, the task is to assign each pixel to one of two disjoint segments. Each pixel carries a unary preference cost reflecting its likelihood of belonging to a specific segment, while each edge carries a pairwise boundary penalty that discourages adjacent pixels from sharing the same segment assignment. The objective is to find a binary labeling of all pixels that minimizes the total segmentation energy, defined as the sum of unary preferences and pairwise boundary penalties.

\paragraph{Decision Variables.}
For each pixel $i \in V$, we introduce a binary decision variable $x_i \in \{0,1\}$. The variable $x_i$ encodes the segment assignment for pixel $i$, where $x_i = 0$ denotes membership in the first segment and $x_i = 1$ denotes membership in the second segment. The complete assignment is represented by the vector $x = (x_i)_{i \in V} \in \{0,1\}^{|V|}$.

\paragraph{QUBO Derivation.}
\textbf{Step 1 --- Original Objective.}
The segmentation energy to be minimized is expressed in general algebraic notation as:
\begin{align}
    \min_{x \in \{0,1\}^{|V|}} \; E(x) = \sum_{i \in V} u_i x_i + \sum_{(i,j) \in E} w_{ij} (1 - x_i x_j),
\end{align}
where $u_i$ denotes the unary preference cost for pixel $i$, and $w_{ij}$ denotes the pairwise boundary penalty between neighboring pixels $i$ and $j$.

\textbf{Step 2 --- Constraints.}
The problem imposes no hard constraints on the assignment. The binary domain $\{0,1\}^{|V|}$ inherently satisfies the requirement that each pixel receives exactly one of two segment labels.

\textbf{Step 3 --- Penalty Terms.}
Since there are no constraints to enforce, no penalty terms are constructed. The formulation relies entirely on the unconstrained objective function.

\textbf{Step 4 --- Combined QUBO Hamiltonian.}
Expanding the objective in Step 1 and collecting terms yields the QUBO Hamiltonian $H(x)$:
\begin{align}
    H(x) = \sum_{i \in V} u_i x_i + \sum_{(i,j) \in E} w_{ij} - \sum_{(i,j) \in E} w_{ij} x_i x_j.
\end{align}
This expression is already in the standard QUBO form $\sum_{i \in V} Q_{ii} x_i + \sum_{i < j} Q_{ij} x_i x_j + c$.

\textbf{Step 5 --- Q Matrix Entries and Offset.}
By direct comparison with the standard QUBO form, the diagonal entries, off-diagonal entries, and constant offset are read off in closed form as:
\begin{align}
    Q_{ii} &= u_i, & \forall i \in V, \\
    Q_{ij} &= -w_{ij}, & \forall (i,j) \in E, \\
    Q_{ij} &= 0, & \forall (i,j) \notin E, \\
    c &= \sum_{(i,j) \in E} w_{ij}.
\end{align}
The coefficients $Q_{ii}$ directly encode the unary pixel-label preference costs, while the coefficients $Q_{ij}$ encode the pairwise boundary penalties between neighboring pixels.

\paragraph{Penalty Weight Justification.}
The problem formulation contains no hard constraints; the binary domain $\{0,1\}^{|V|}$ naturally enforces the two-segment assignment requirement for every pixel. Consequently, no penalty weights $\lambda$ are required, and the QUBO mapping is exact without the need for constraint relaxation or penalty scaling.

\paragraph{Correctness Argument.}
Minimizing $H(x)$ over $\{0,1\}^{|V|}$ recovers the optimal segmentation because the feasible set of the QUBO exactly matches the feasible set of the original problem. Every binary assignment corresponds to a valid two-segment partition, and since no penalty terms are added, the energy of any assignment is identical to its original segmentation energy. Therefore, the global minimum of $H(x)$ must occur at a feasible assignment that minimizes the original objective, guaranteeing that the optimal $x^*$ returned by the solver corresponds to the minimum-energy segmentation.

\paragraph{Illustrative Example.}
Consider a concrete instance with $n = 4$ pixels arranged in a $2 \times 2$ grid. The vertex set is $V = \{0, 1, 2, 3\}$ and the edge set is $E = \{(0,1), (0,2), (1,3), (2,3)\}$. We instantiate the general formulas with specific cost parameters: $u_0 = 1, u_1 = 2, u_2 = 3, u_3 = 4$ and $w_{01} = 1, w_{02} = 2, w_{13} = 3, w_{23} = 4$. Substituting these into the general Hamiltonian yields:
\begin{align}
    H(x) &= 1x_0 + 2x_1 + 3x_2 + 4x_3 + (1+2+3+4) - (1x_0x_1 + 2x_0x_2 + 3x_1x_3 + 4x_2x_3) \\
         &= 1x_0 + 2x_1 + 3x_2 + 4x_3 - 1x_0x_1 - 2x_0x_2 - 3x_1x_3 - 4x_2x_3 + 10.
\end{align}
The corresponding Q matrix entries and offset are:
\begin{align}
    Q_{00} &= 1, & Q_{11} &= 2, & Q_{22} &= 3, & Q_{33} &= 4, \\
    Q_{01} &= -1, & Q_{02} &= -2, & Q_{13} &= -3, & Q_{23} &= -4, \\
    c &= 10.
\end{align}
Evaluating $H(x)$ over all $2^4 = 16$ binary assignments reveals that the minimum energy is $E_{\min} = 10$, achieved at $x^* = (0, 0, 0, 0)$. This assignment corresponds to placing all four pixels in the first segment, which is optimal given the relatively low unary costs and the absence of boundary penalties for uniform labeling in this specific parameterization. For any other assignment, the increase in unary costs or the reduction of negative pairwise contributions results in strictly higher energy.
\end{document}